% Capitulo 2: flujo de trabajo y comandos, con la tabla de comandos.
\section{El flujo de trabajo}\label{sec:flujo}

Un proyecto texforge es un directorio con un \texttt{project.toml} que nombra
la entrada y la bibliografía. Se empieza con \texttt{texforge new
  mi-documento}, que genera el esqueleto, o con \texttt{texforge init} dentro
de un proyecto \LaTeX{} existente, que detecta el \texttt{main.tex} y la
bibliografía y escribe el \texttt{project.toml} por ti. A partir de ahí
el ciclo es corto: escribir, revisar con \texttt{check}, formatear con
\texttt{fmt} y componer con \texttt{build}. \texttt{texforge clean} retira el
directorio \texttt{build/} cuando quieres empezar la composición desde
cero, y \texttt{texforge outline} muestra el árbol de secciones del
documento para comprobar la estructura sin abrir el PDF.

\subsection{Inspeccionar el resultado}

Componer no es el último paso: el PDF también se lee desde la terminal.
\texttt{texforge pdf text} vuelca el texto tal como lo vería un lector de
pantalla o un sistema de selección de personal, con las ligaduras
normalizadas y las palabras partidas al final de línea ya reunidas, así
que es la prueba real de accesibilidad del documento. \texttt{texforge pdf
  info} informa de las páginas, de las fuentes y de si van embebidas, y de
los metadatos que viajarán con el archivo. \texttt{texforge pdf pages}
dice qué sección abre cada página en un formato pensado para comparar
compilaciones línea a línea, y \texttt{texforge pdf check} comprueba que
cada palabra significativa del fuente haya sobrevivido a la composición:
si una ligadura o una codificación se come una palabra, el aviso aparece
aquí y no en la imprenta. Para mirar el resultado con los ojos,
\texttt{texforge preview} convierte el PDF en imágenes sin salir de la
terminal. Y cuando el corrector duda de una palabra, \texttt{texforge spell}
decide: las palabras personales, como apellidos o tecnologías, se aceptan
una vez para todos los proyectos, y las de un documento concreto se guardan
con \texttt{-{}-local} en el propio proyecto, en \texttt{.texforge/spell-words}.
La lista de este documento, sin ir más lejos, guarda los nombres de las
herramientas y los términos técnicos que el diccionario no trae, para que
\texttt{check} revise prosa y no vocabulario.

\subsection{Los comandos}

La tabla~\ref{tab:comandos} resume los comandos, incluidos los que llegaron
con la versión 1.0. Las novedades de la versión se describen con calma en
el capítulo~\ref{sec:novedades}; aquí basta saber que el flujo diario
cabe en \texttt{new}, \texttt{build}, \texttt{check} y \texttt{fmt}, y que el
resto aparece cuando el documento crece: plantillas, ortografía,
diagnóstico y lectura del PDF.

\begin{table}[H]
  \centering
  \caption{Comandos de texforge, con las novedades de la versión 1.0}
  \label{tab:comandos}
  \begin{tabularx}{\linewidth}{@{}lX@{}}
    \toprule
    \sffamily\small\bfseries Comando & \sffamily\small\bfseries Qué hace \\
    \midrule
    \texttt{new -t -{}-values} & Crea un proyecto rellenando la plantilla desde un fichero de valores, que queda guardado como \texttt{values.toml}. \\
    \texttt{template add} & Añade una carpeta local como plantilla, con \texttt{-{}-force} para reemplazarla. \\
    \texttt{skill} & Muestra la destreza de agente que viaja dentro del binario. \\
    \texttt{doctor} & Diagnostica el entorno y actualiza el Tectonic gestionado; avisa si el de \texttt{PATH} es más viejo. \\
    \texttt{update} & Actualiza texforge verificando la firma minisign del paquete. \\
    \texttt{outline} & Muestra el árbol de secciones del documento. \\
    \texttt{pdf text} & Imprime el texto del PDF como lo vería un lector o un sistema de selección. \\
    \texttt{pdf info} & Informa de páginas, fuentes embebidas y metadatos. \\
    \texttt{pdf pages} & Dice qué sección abre cada página, en un formato apto para comparar compilaciones. \\
    \texttt{pdf check} & Comprueba que las palabras del fuente sobrevivieron a la composición. \\
    \texttt{build} & Compone \texttt{main.tex} a PDF en \texttt{build/}. \\
    \texttt{check} & Revisa citas, etiquetas y ortografía sin compilar. \\
    \texttt{fmt} & Formatea los \texttt{.tex} y los \texttt{.bib}. \\
    \texttt{init} & Migra un proyecto \LaTeX{} existente o guía uno nuevo. \\
    \texttt{clean} & Retira el directorio \texttt{build/}. \\
    \texttt{stats} & Cuenta palabras por sección o por archivo. \\
    \texttt{preview} & Convierte el PDF en imágenes para revisarlo. \\
    \texttt{config} & Lee y escribe la configuración global. \\
    \bottomrule
  \end{tabularx}
\end{table}

\begin{novabox}
  \texttt{outline} y los cuatro subcomandos de \texttt{pdf} son nuevos en la
  1.0: antes la estructura y el PDF resultante solo se veían abriendo el
  PDF; ahora se inspeccionan desde la terminal.
\end{novabox}

Revisar antes de componer ahorra tiempo. \texttt{texforge check} encuentra
archivos \texttt{\textbackslash input} que faltan, imágenes ausentes,
citas sin entrada en el \texttt{.bib} y pares \texttt{\textbackslash ref} y
\texttt{\textbackslash label} rotos, incluidos los \texttt{label=} de los
diagramas. El corrector ortográfico respeta el idioma declarado con
\texttt{babel} y acepta el vocabulario propio con \texttt{texforge spell}:
las palabras personales viven en el diccionario global y las de un documento
concreto se guardan con \texttt{-{}-local} en el proyecto. \texttt{texforge
  fmt} deja la indentación y los \texttt{.bib} en forma canónica, y
\texttt{fmt -{}-check} lo verifica en integración continua sin tocar nada.
